\frontmatter
% v1.1 AI statement. Preserve the 12-column form from ThesisPlan PDF p.6.
% Requires pdflscape, longtable, array, graphicx, multirow, amssymb and setspace.
\chapter*{Use of AI Statement}
\addcontentsline{toc}{chapter}{Use of AI Statement}
\noindent\textbf{v1.1}

I acknowledge the use of AI tools in this project. The table records their use from the first-semester proposal through the present thesis. I take responsibility for evaluating and integrating the assisted material and for the final submitted work.

I conducted the study and formed its findings and key observations. I supplied the behavioural accounts, selected the important cases, and corrected their interpretation. AI assisted with checking those accounts against selected source records, organising the material, computing and drawing selected representations, and expressing the work in English. It also supported earlier planning, literature reading, Unity prototype coding and debugging, research-logging implementation, and document preparation.

The original Thesis Plan table structure is retained below. A tick identifies a documented use in that stage. ``Data analysis'' includes assisted data organisation, numerical checks and selected log or geometric verification; it does not attribute the discovery of the findings to AI. ``Visual content creation'' includes document formatting and programmatic figures made from recorded data. No participant observations or questionnaire answers were generated by AI. Software coding, debugging and data-capture implementation are described under ``Other use'' because the original form has no separate code column. Tool names are reported where exact model versions were not retained. Semester-level dates are used where the early records do not establish a reliable usage day.

\begin{landscape}
\begingroup
\setstretch{1}
\fontsize{10}{12}\selectfont
\setlength{\tabcolsep}{2pt}
\renewcommand{\arraystretch}{1.1}
\setlength{\LTpre}{0pt}
\setlength{\LTpost}{0pt}
\begin{longtable}{|>{\raggedright\arraybackslash}p{0.14\linewidth}|*{10}{>{\centering\arraybackslash}p{0.027\linewidth}|}>{\raggedright\arraybackslash}p{\dimexpr0.59\linewidth-24\tabcolsep-13\arrayrulewidth\relax}|}
\hline
\multirow{2}{*}{\shortstack{\textbf{AI model}\\\textbf{and date}}} & \multicolumn{11}{c|}{\textbf{Artificial Intelligence (AI) tool usage} (tick all that apply)} \\
\cline{2-12}
 & \rotatebox{90}{\strut Topic exploration} & \rotatebox{90}{\strut Research question} & \rotatebox{90}{\strut Literature review} & \rotatebox{90}{\strut Methodological support} & \rotatebox{90}{\strut Data analysis} & \rotatebox{90}{\strut Visual content creation} & \rotatebox{90}{\strut Written content creation} & \rotatebox{90}{\strut Language translation} & \rotatebox{90}{\strut Grammar/style/spelling} & \rotatebox{90}{\strut Feedback on content} & \multicolumn{1}{c|}{\shortstack{Other use\\(provide details)}} \\
\hline
\endfirsthead
\hline
\multirow{2}{*}{\shortstack{\textbf{AI model}\\\textbf{and date}}} & \multicolumn{11}{c|}{\textbf{Artificial Intelligence (AI) tool usage} (tick all that apply)} \\
\cline{2-12}
 & \rotatebox{90}{\strut Topic exploration} & \rotatebox{90}{\strut Research question} & \rotatebox{90}{\strut Literature review} & \rotatebox{90}{\strut Methodological support} & \rotatebox{90}{\strut Data analysis} & \rotatebox{90}{\strut Visual content creation} & \rotatebox{90}{\strut Written content creation} & \rotatebox{90}{\strut Language translation} & \rotatebox{90}{\strut Grammar/style/spelling} & \rotatebox{90}{\strut Feedback on content} & \multicolumn{1}{c|}{\shortstack{Other use\\(provide details)}} \\
\hline
\endhead
\hline
\multicolumn{12}{r}{\textit{Continued on the next page}}\\
\endfoot
\hline
\endlastfoot
\textbf{ChatGPT}\newline Semester 1, 2026 & \ensuremath{\checkmark} & \ensuremath{\checkmark} &  & \ensuremath{\checkmark} &  &  & \ensuremath{\checkmark} & \ensuremath{\checkmark} & \ensuremath{\checkmark} & \ensuremath{\checkmark} & \textbf{Topic definition and early study planning.} Helped organise the author's ideas, discuss possible research framings and study structures, and turn early planning notes into clearer text. The author and supervisor made the research decisions. \\
\hline
\textbf{ChatGPT}\newline Semester 1, 2026 &  &  & \ensuremath{\checkmark} &  &  &  & \ensuremath{\checkmark} & \ensuremath{\checkmark} & \ensuremath{\checkmark} & \ensuremath{\checkmark} & \textbf{Literature reading and proposal writing.} Explained difficult passages, compared the relevance of papers located by the author, and assisted with proposal organisation, Chinese-to-English translation and wording. Literature claims remained tied to the cited sources. \\
\hline
\textbf{Gemini Pro}\newline Semester 1, 2026 &  &  &  &  &  & \ensuremath{\checkmark} & \ensuremath{\checkmark} & \ensuremath{\checkmark} & \ensuremath{\checkmark} &  & \textbf{Proposal typesetting.} Assisted with translation, language correction, visual formatting and LaTeX code for the proposal, as disclosed in that document. Visual content creation here refers to document formatting. \\
\hline
\textbf{ChatGPT}\newline Semester 1, 2026 &  &  &  & \ensuremath{\checkmark} &  &  & \ensuremath{\checkmark} & \ensuremath{\checkmark} & \ensuremath{\checkmark} & \ensuremath{\checkmark} & \textbf{Ethics wording and meeting notes.} Assisted with wording for draft ethics and study-planning documents, and with organising supervisor-meeting notes and follow-up tasks. \\
\hline
\textbf{ChatGPT / OpenAI Codex}\newline 17--19 August 2026 & \ensuremath{\checkmark} & \ensuremath{\checkmark} & \ensuremath{\checkmark} & \ensuremath{\checkmark} &  &  & \ensuremath{\checkmark} & \ensuremath{\checkmark} & \ensuremath{\checkmark} & \ensuremath{\checkmark} & \textbf{Thesis Plan.} Reorganised existing proposal and project material for the Thesis Plan, compared research framings, drafted connecting English text, and prepared LaTeX formatting. \\
\hline
\textbf{OpenAI Codex}\newline September 2026 &  & \ensuremath{\checkmark} &  & \ensuremath{\checkmark} &  &  & \ensuremath{\checkmark} & \ensuremath{\checkmark} & \ensuremath{\checkmark} & \ensuremath{\checkmark} & \textbf{Method and prototype preparation.} Reviewed the author's plans and supervisor feedback, helped refine questionnaire and interview procedures, and organised prototype, viewing-orientation and recording requirements. The author supplied revisions and selected the implemented procedure. \\
\hline
\textbf{OpenAI Codex}\newline 20--23 September 2026 &  &  &  & \ensuremath{\checkmark} &  &  & \ensuremath{\checkmark} &  &  &  & \textbf{Multiplayer prototype coding and debugging.} Wrote and revised Unity C\# scripts and scene configuration for the two-player prototype, object synchronisation, grab/release control, first-person and overhead display, connection menus and builds. Debugging followed the author's headset-test feedback; technical instructions were also prepared. \\
\hline
\textbf{OpenAI Codex}\newline 6--7 October 2026 &  &  &  & \ensuremath{\checkmark} &  &  & \ensuremath{\checkmark} &  &  &  & \textbf{Research logging implementation.} Implemented and checked the ResearchV2 CSV logger and its connections to piece actions, recorded poses, round controls and participant/round identifiers. Prepared validation code, builds and export guidance. \\
\hline
\textbf{OpenAI Codex}\newline 6--7 October 2026 &  &  &  & \ensuremath{\checkmark} &  &  & \ensuremath{\checkmark} & \ensuremath{\checkmark} & \ensuremath{\checkmark} & \ensuremath{\checkmark} & \textbf{Collection materials.} Revised the Chinese collection guide, moderator manual, researcher record sheet and twelve-item questionnaire; documented Unity log fields and later evidence use. \\
\hline
\textbf{OpenAI Codex}\newline 7--10 October 2026 &  &  &  & \ensuremath{\checkmark} &  &  & \ensuremath{\checkmark} &  &  &  & \textbf{Prototype revisions and technical checks.} Implemented the shared target, manual round ending and overhead-menu repair, then target grabbing, ray interaction and the later view adjustment. Produced checks and APK builds from the author's requirements and reported faults. Each revision retained its version identity. \\
\hline
\textbf{ChatGPT / OpenAI Codex}\newline 13--14 October 2026 &  &  & \ensuremath{\checkmark} &  & \ensuremath{\checkmark} &  & \ensuremath{\checkmark} & \ensuremath{\checkmark} & \ensuremath{\checkmark} & \ensuremath{\checkmark} & \textbf{Conference paper.} Organised researcher-provided case records and questionnaire summaries, supported comparison of evidence across sources, and helped draft, translate, edit and format the conference paper. The author supplied observations and participant responses. \\
\hline
\textbf{ChatGPT / OpenAI Codex}\newline 16--21 October 2026 &  &  &  & \ensuremath{\checkmark} & \ensuremath{\checkmark} & \ensuremath{\checkmark} & \ensuremath{\checkmark} &  & \ensuremath{\checkmark} & \ensuremath{\checkmark} & \textbf{Corpus organisation and data reconciliation.} Indexed videos, reports, interviews and logs; organised the author's observations; checked selected records; reconciled the 48 confirmed questionnaire responses; and produced descriptive tables, evidence maps and programmatic plots. \\
\hline
\textbf{TurboScribe}\newline Records retrieved 20 October 2026 &  &  &  &  &  &  &  &  &  &  & \textbf{Meeting transcription.} Automatic transcription of recorded supervisor meetings used to recover project history. Existing text and subtitle exports were retrieved on 20 October; this is an archive date, not a claimed transcription date for every recording. \\
\hline
\textbf{ChatGPT / OpenAI Codex}\newline 21--22 October 2026 &  &  & \ensuremath{\checkmark} &  & \ensuremath{\checkmark} & \ensuremath{\checkmark} & \ensuremath{\checkmark} & \ensuremath{\checkmark} & \ensuremath{\checkmark} & \ensuremath{\checkmark} & \textbf{Poster and presentation.} Checked selected action sequences and reconstructed geometry for the author's observations; prepared programmatic figures, poster text and layout; and revised speaking notes from the author's oral explanations. The author chose the findings and corrected the accounts. \\
\hline
\textbf{ChatGPT / OpenAI Codex}\newline 25 October--2 November 2026 &  &  & \ensuremath{\checkmark} & \ensuremath{\checkmark} & \ensuremath{\checkmark} & \ensuremath{\checkmark} & \ensuremath{\checkmark} & \ensuremath{\checkmark} & \ensuremath{\checkmark} & \ensuremath{\checkmark} & \textbf{Thesis integration and revision.} Organised thesis sources and sections; checked selected logs, geometry and references; produced code-based figures; drafted and revised English from the author's material; and prepared LaTeX, tables, this disclosure and PDF checks. \\
\hline
\end{longtable}
\endgroup
\end{landscape}

\clearpage
\chapter*{Abstract}
\addcontentsline{toc}{chapter}{Abstract}

A shared overhead view gives collaborators access to the same spatial arrangement while they act on a developing construction. This thesis examines how that view relates to shape judgment and the coordination of work in collaborative Tangram solving. Six pairs took part in four conditions: physical direct viewing, physical overhead video, first-person VR and shared overhead VR. The analysis follows their construction processes through the researcher's observations, video recordings, selected action and pose logs, questionnaires and interviews.

Three connected findings emerged. First, near-complete constructions in first-person VR could look convincing while leaving a piece unused or differing from the target outline. The same class of problem was not observed in the shared-overhead VR records, where the planar relationships were directly visible. Second, participants expressed ideas by placing pieces, and their partners could continue those placements or demonstrate an alternative. The division of work changed during a task as people moved between separate attempts, assistance and joint reconstruction. Third, several episodes showed a partner retaining or extending a local construction that its original builder later revised, including cases where a correct part was lost. Other episodes showed direct partner intervention when an alternative was put into action.

These findings connect the geometric advantage of an overhead view with the actions through which partners develop and judge a shared solution. They support combining an inspectable overview with usable manipulation, brief assistance and access to earlier constructions during revision.

\vspace{0.6em}
\noindent\textbf{Keywords:} collaborative virtual reality; shared overhead view; Tangram; shape judgment; action coordination

\clearpage
\tableofcontents
\clearpage
\listoffigures
\clearpage
\listoftables
